\scanpage{38}
% archive sheet number: 37

Soit \(G\) un schéma en groupes simple sur \(S=\operatorname{Spec} A\), \(A\) Artinien local.

\[
\mathfrak g=L_S(G)
\]

(1) Si \(H^0(\mathfrak g_0,\mathfrak g_0)=0\), \deleted{\unsure{alors}}\(\Longleftrightarrow\) \(G\longrightarrow\underline{\operatorname{Aut}}_S(\mathfrak g)\) est \deleted{\unsure{\illegible{1}|un}} non ramifié

[donc si \(G_0\) n'a pas de sous-groupe \deleted{discret} étale \deleted{\unsure{forcément|forcément un}} non trivial, alors c'est une imm. fermée]

(2) \deleted{Si}\inserted{(1) et} \inserted{\deleted{de plus}} \(H^1(\mathfrak g_0,\mathfrak g_0)=0\) \deleted{\unsure{alors}}\(\Longleftrightarrow\) \(G\longrightarrow\underline{\operatorname{Aut}}_S(\mathfrak g)\) est étale,

[donc sous la condition dite plus haut, c'est une imm. ouverte, donc \deleted{\unsure{Aut}}\(G\simeq\underline{\operatorname{Aut}}_S^{\circ}(\mathfrak g)\) \deleted{\unsure{G}}si \(G\) à fibres connexes]

% slanted
\margin{\deleted{(1) et (2)} (2) \(\Longrightarrow\) \(\underline{\operatorname{Aut}}_S(\mathfrak g)\simeq\underline{\operatorname{Aut}}_S(G)\) si \(G\) connexe \ldots\ Utilisant \unsure{Tchoka|Tôhoku|Tohoku}, cela suffirait à construire les \unsure{Tchokas|Tôhokus|Tohoku} \deleted{\illegible{1}} sur \(Z\) \ldots}

Supposons qu'on ne dispose pas de \(G\), mais seulement de \(\mathfrak g\), considérons \(H=\underline{\operatorname{Aut}}_S(\mathfrak g)\). \deleted{Alors} C'est un schém. en groupes sur \(S\), et \(\operatorname{Lie}_k(H_0)=\mathfrak h_0\) est isomorphe à \deleted{\(\mathfrak g_0/(\text{centre de }\mathfrak g_0)\) par \unsure{iso. canonique|son caractère|iso}} \inserted{l'Algèbre de Lie des} dérivations de \(\mathfrak g_0\), d'où \(\mathfrak g_0\to\mathfrak h_0\), qui est un isom. si (1) et (2) sont satisfaites.

S'il en est ainsi, on ne peut affirmer encore que \deleted{\unsure{G}}\(H\) est simple sur \(S\), \unsure{\(=\)|\(\simeq\)} le long de la section unité, \deleted{Mais dans} \unsure{On|Or} le peut si on suppose

\[
\text{(3)}\qquad H^2(\mathfrak g_0,\mathfrak g_0)=0
\]

qui signifie sauf erreur qu'il n'y a pas de \unsure{variation|variations} de structure \inserted{\unsure{infinitésimales|infinitésimale}} non triviales de \(\mathfrak g_0\), et \uline{implique} que \deleted{\unsure{les deux|tous deux}} \(\underline{\operatorname{Aut}}_S(\mathfrak g)\) est simple sur \(S\).
